% Created 2026-02-09 seg 15:06
% Intended LaTeX compiler: lualatex
\documentclass[12pt]{article}
\usepackage[a4paper,left=3cm,top=3cm,right=2cm,bottom=2cm]{geometry}
\setcounter{secnumdepth}{4}
\setcounter{tocdepth}{2}
\usepackage{graphicx}
\usepackage{longtable}
\usepackage{wrapfig}
\usepackage{rotating}
\usepackage[normalem]{ulem}
\usepackage{amsmath}
\usepackage{amssymb}
\usepackage{capt-of}
\usepackage{hyperref}
\usepackage{fgv-header}
\author{Gustavo M. Mendes de Tarso}
\date{\today}
\title{}
\hypersetup{
 pdfauthor={Gustavo M. Mendes de Tarso},
 pdftitle={},
 pdfkeywords={},
 pdfsubject={},
 pdfcreator={Emacs 28.2 (Org mode 9.5.5)}, 
 pdflang={English}}
\begin{document}

\begingroup
\linespread{1}\selectfont
\begin{tcolorbox}[title=Ficha Técnica,
  colback=gray!5,colframe=gray!40,boxrule=0.4pt,sharp corners]
\begin{tblr}{rowsep=1pt,stretch=1, rows={t}, colspec={Q[l,2.8cm] X[l]}}
\textbf{Curso}   		& MBA em Gestão da Saúde \\
\textbf{Turma}      		& T-01 \\
\textbf{Pólo}      		& Brasília \\
\textbf{Disciplina}     	& Gestão de Custos \\
\textbf{Professor}      	& Waldir Jorge Ladeira dos Santos \\
\textbf{Trabalho}  	        & Resolução de exercício da apostila \\
\textbf{Aluno(s)}     		& Gustavo M. Mendes de Tarso. \\
\textbf{Prazo}     		& 09 de fevereiro de 2026 \\
\textbf{Título do trabalho}     & Gestão de Custos do Hospital Alfa \\ 
\end{tblr}
\end{tcolorbox}
\endgroup

\FGVNewPage

\section{Exercício 4 — Hospital Alfa (Novembro)}
\label{sec:org9754076}

\subsection{Premissas do enunciado}
\label{sec:org618c7c5}
\begin{itemize}
\item Receita (serviços prestados) no período: \textbf{R\$ 2.100.000}.
\item Cenário (b): \textbf{vendas +10\%}.
\item Custos e despesas variáveis variam na mesma proporção das vendas.
\item Custos e despesas fixas permanecem inalteradas.
\end{itemize}

\subsection{Classificação dos itens (Custo/Despesa e Variável/Fixo)}
\label{sec:org0c48861}

\begin{table}[htbp]
\centering
\scriptsize
\setlength{\tabcolsep}{3pt}
\renewcommand{\arraystretch}{1.15}
\begin{tabularx}{\textwidth}{%
>{\raggedright\arraybackslash}X
>{\raggedleft\arraybackslash}m{0.14\textwidth}
>{\centering\arraybackslash}m{0.18\textwidth}
>{\centering\arraybackslash}m{0.18\textwidth}}
\toprule
Itens & Valor (R\$) & DRE Contábil/Financeira & DRE Gerencial/Interna \\
\midrule
Depreciação dos equipamentos hospitalares & 80.000 & Custo & Fixo \\
Custo de expedição interna da farmácia & 25.000 & Custo & Variável \\
Material hospitalar diretamente aplicado & 350.000 & Custo & Variável \\
Mão-de-obra direta (médico, enfermeiro...) & 400.000 & Custo & Variável \\
Propaganda contratada de forma variável & 58.000 & Despesa (Comercial) & Variável \\
Funcionários da cozinha para os enfermos & 57.000 & Custo & Fixo \\
Seguro contra-incêndio do hospital (área fim) & 15.000 & Custo & Fixo \\
Imposto sobre propriedade do hospital (área fim) & 25.100 & Custo & Fixo \\
Energia, água e esgoto dos procedimentos & 59.500 & Custo & Variável \\
Ordenados do pessoal de Marketing & 25.500 & Despesa (Comercial) & Fixo \\
Despesas administrativas & 50.000 & Despesa (Administrativa) & Fixo \\
Aluguel do hospital (área fim) & 90.000 & Custo & Fixo \\
\bottomrule
\end{tabularx}
\normalsize
\end{table}

\subsection{DRE — Abordagem Funcional / Absorção}
\label{sec:org733ee7a}

\begin{table}[htbp]
\centering
\scriptsize
\setlength{\tabcolsep}{3pt}
\renewcommand{\arraystretch}{1.15}
\begin{tabularx}{\textwidth}{%
>{\raggedright\arraybackslash}X
>{\raggedleft\arraybackslash}m{0.20\textwidth}
>{\raggedleft\arraybackslash}m{0.20\textwidth}}
\toprule
DRE Abordagem Funcional (Absorção) & a) & b) (vendas +10\%) \\
\midrule
Receita de venda & 2.100.000 & 2.310.000 \\
(--) Custos dos serviços prestados & 1.101.600 & 1.185.050 \\
\midrule
\textbf{Lucro bruto} & \textbf{998.400} & \textbf{1.124.950} \\
\midrule
(--) Despesas operacionais (total) & 133.500 & 139.300 \\
\hspace{2mm} Administrativas & 50.000 & 50.000 \\
\hspace{2mm} Comerciais & 83.500 & 89.300 \\
\midrule
\textbf{Lucro operacional líquido} & \textbf{864.900} & \textbf{985.650} \\
\bottomrule
\end{tabularx}
\normalsize
\end{table}

\subsection{DRE — Abordagem de Contribuição / Variável}
\label{sec:org21eee2e}

\begin{table}[htbp]
\centering
\scriptsize
\setlength{\tabcolsep}{3pt}
\renewcommand{\arraystretch}{1.15}
\begin{tabularx}{\textwidth}{%
>{\raggedright\arraybackslash}X
>{\raggedleft\arraybackslash}m{0.20\textwidth}
>{\raggedleft\arraybackslash}m{0.20\textwidth}}
\toprule
DRE Abordagem de Contribuição (Variável) & a) & b) (vendas +10\%) \\
\midrule
Receita de venda & 2.100.000 & 2.310.000 \\
(--) Custos variáveis & 834.500 & 917.950 \\
(--) Despesas variáveis & 58.000 & 63.800 \\
\midrule
\textbf{Margem de contribuição} & \textbf{1.207.500} & \textbf{1.328.250} \\
\midrule
(--) Custos fixos & 267.100 & 267.100 \\
(--) Despesas fixas & 75.500 & 75.500 \\
\midrule
\textbf{Lucro operacional líquido} & \textbf{864.900} & \textbf{985.650} \\
\bottomrule
\end{tabularx}
\normalsize
\end{table}

\subsection{Qual demonstrativo usar no item (b) e por quê?}
\label{sec:org16a58fb}
A DRE mais adequada para o item (b) é a \textbf{DRE pela Abordagem de Contribuição (Variável)}, pois o
enunciado estabelece explicitamente que custos e despesas variáveis variam na mesma proporção
das vendas e que custos e despesas fixas permanecem inalteradas. Como esse demonstrativo
separa de forma objetiva os componentes variáveis dos fixos, ele permite mensurar diretamente
o impacto do aumento de 10\% na receita sobre a margem de contribuição e, por consequência,
sobre o lucro operacional, sendo portanto o modelo gerencial mais apropriado para a simulação
custo-volume-lucro.
\end{document}